% block: jordan-canonical-heading page:201 type:section
\subsection{{Jordan 标准形}}\label{sec:jordan-canonical-form}

% block: jordan-canonical-choice page:201 type:ordinary-paragraph
{在矩阵函数的若干等价定义中, 我们采用利用 Jordan 标准形的定义. 本节介绍 Jordan 标准形.}

% block: jordan-block-definition pages:201,202 type:definition
\Needspace{8\baselineskip}
\begin{Definition}{{Jordan 块}}\label{def:jordan-block}
{下面的 $m\times m$ 方阵称为 Jordan 块:}
\[
J_k=\begin{bmatrix}\lambda_k&1&&\\&\lambda_k&\ddots&\\&&\ddots&1\\&&&\lambda_k\end{bmatrix}\in\mathbb{C}^{m\times m},
\]
{其中 $\lambda_k\in\mathbb{C}$. 由于 $J_k$ 是上三角矩阵, $J_k$ 的所有特征值都与对角元素相同, 即 $\lambda_k$.}
\end{Definition}

% block: jordan-matrix-definition page:202 type:definition
\Needspace{9\baselineskip}
\begin{Definition}{{Jordan 矩阵}}\label{def:jordan-matrix}
{Jordan 块 $J_1\in\mathbb{C}^{m_1\times m_1}$, $J_2\in\mathbb{C}^{m_2\times m_2}$, $\ldots$, $J_p\in\mathbb{C}^{m_p\times m_p}$ 的直和为}
\begin{equation}\label{eq:jordan-block-diagonal}
J=\operatorname{diag}(J_1,J_2,\ldots,J_p)=\begin{bmatrix}J_1&&&\\&J_2&&\\&&\ddots&\\&&&J_p\end{bmatrix}\in\mathbb{C}^{n\times n}.
\end{equation}
{该矩阵称为 Jordan 矩阵. 这里, $n=m_1+m_2+\cdots+m_p$, 符号 $\operatorname{diag}(J_1,J_2,\ldots,J_p)$ 表示以 $J_1,J_2,\ldots,J_p$ 为对角块的矩阵.}
\end{Definition}

% block: jordan-six-dimensional-example page:202 type:example
\Needspace{10\baselineskip}
\begin{Example}\label{ex:jordan-block-diagonal-six}
{设 $p=3$, $m_1=2$, $m_2=1$, $m_3=3$, $\lambda_1=1$, $\lambda_2=\lambda_3=2$. 则 $n=m_1+m_2+m_3=6$, 并得到下面的 $6\times6$ Jordan 矩阵:}
\[
J=\operatorname{diag}(J_1,J_2,J_3)=\left[\begin{array}{cc|c|ccc}1&1&0&0&0&0\\0&1&0&0&0&0\\\hline0&0&1&0&0&0\\\hline0&0&0&2&1&0\\0&0&0&0&2&1\\0&0&0&0&0&2\end{array}\right].
\]
\end{Example}

% block: jordan-canonical-theorem page:202 type:theorem
\Needspace{14\baselineskip}
\begin{Theorem}{{Jordan 标准形}}\label{thm:jordan-canonical-form}
{方阵 $A\in\mathbb{C}^{n\times n}$ 相似于一个 Jordan 矩阵, 即对任意方阵 $A$, 存在非奇异矩阵 $Z$ 使得}
\[Z^{-1}AZ=J.\]
{形式 $A=ZJZ^{-1}$ 称为 $A$ 的 Jordan 标准形.}
\end{Theorem}

% block: jordan-similarity-explicit-example pages:202,203 type:example
\Needspace{10\baselineskip}
\begin{Example}\label{ex:jordan-similarity-six}
{下面的矩阵相似于 Example~\ref{ex:jordan-block-diagonal-six} 中的 Jordan 矩阵 $J$.}
\[
A=\begin{bmatrix}-1&1&-1&-2&-2&-1\\1&2&2&1&2&1\\-3&-3&-4&-3&-5&-3\\4&0&3&5&4&2\\4&4&7&4&8&4\\-3&-2&-5&-3&-5&-1\end{bmatrix}.
\]
{事实上, 可以看到 $Z^{-1}AZ=J$, 其中}
\[
Z=\begin{bmatrix}1&0&0&-1&0&0\\0&1&-1&0&1&-1\\0&-1&2&0&-2&1\\-1&0&0&2&0&-1\\0&1&-2&0&3&-2\\0&-1&1&-1&-2&4\end{bmatrix},\quad
Z^{-1}=\begin{bmatrix}3&1&1&2&1&1\\1&3&2&1&1&1\\1&2&3&1&2&1\\2&1&1&2&1&1\\1&1&2&1&2&1\\1&1&1&1&1&1\end{bmatrix}.
\]
\end{Example}

% block: jordan-definition-transition page:203 type:ordinary-paragraph
{\leavevmode\phantomsection\label{srcpage:printed-191}在矩阵函数的几种定义中, 下一节采用利用 Theorem~\ref{thm:jordan-canonical-form} 中 Jordan 标准形的定义.}
